% !TeX program = xelatex
% !BIB program = biber
% ============================================================================
% TFM — Máster en Finanzas Cuantitativas (UNED, Formación Permanente)
% Compilación recomendada en VS Code / LaTeX Workshop:
%   latexmk -xelatex main.tex
% Si se añade bibliografía:
%   xelatex main.tex -> biber main -> xelatex main.tex -> xelatex main.tex
% ============================================================================

\documentclass[12pt,a4paper,oneside,openany]{report}

% ----------------------------------------------------------------------------
% IDIOMA Y TIPOGRAFÍA
% ----------------------------------------------------------------------------
\usepackage[english,spanish]{babel}
\usepackage{fontspec}
\defaultfontfeatures{Ligatures=TeX}
\usepackage{microtype}

% En Windows se utilizará Times New Roman. El fallback permite compilar el
% proyecto en equipos donde esa fuente no esté instalada.
\IfFontExistsTF{Times New Roman}
  {\setmainfont{Times New Roman}}
  {\setmainfont{TeX Gyre Termes}}

% Matemáticas con estética compatible con Times.
\usepackage{amsmath,amssymb,mathtools}
% newtx needs a legacy Times family for math alphabets, not the TU fontspec name.
\let\TFMTextFamily\rmdefault
\renewcommand{\rmdefault}{ptm}
\usepackage{newtxmath}
\let\rmdefault\TFMTextFamily

% ----------------------------------------------------------------------------
% PÁGINA, PÁRRAFOS Y ESPACIADO
% ----------------------------------------------------------------------------
\usepackage[
  a4paper,
  top=2.5cm,
  bottom=2.5cm,
  left=2.5cm,
  right=2.5cm,
  footskip=1.2cm
]{geometry}

\usepackage{setspace}
\onehalfspacing

% Estilo inspirado en el TFG previo: texto justificado, sangría de primera
% línea y sin espacio vertical adicional entre párrafos. LaTeX no sangra el
% primer párrafo tras un título, que es el comportamiento deseado.
\setlength{\parindent}{1.25cm}
\setlength{\parskip}{0pt}
\setlength{\emergencystretch}{2em}
\raggedbottom

% Evita grandes huecos verticales sin comprimir en exceso el texto.
\setlength{\abovedisplayskip}{8pt plus 2pt minus 3pt}
\setlength{\belowdisplayskip}{8pt plus 2pt minus 3pt}
\setlength{\abovedisplayshortskip}{5pt plus 2pt minus 2pt}
\setlength{\belowdisplayshortskip}{5pt plus 2pt minus 2pt}
\allowdisplaybreaks[2]

% Control moderado de viudas y huérfanas.
\widowpenalty=7000
\clubpenalty=7000
\displaywidowpenalty=7000

% ----------------------------------------------------------------------------
% JERARQUÍA DE TÍTULOS
% ----------------------------------------------------------------------------
\usepackage{titlesec}

% Capítulos: 14 pt, negrita, centrados, numeración "1. Título".
\titleformat{\chapter}[block]
  {\normalfont\bfseries\centering\fontsize{14}{17}\selectfont}
  {\thechapter.}
  {0.6em}
  {}
\titlespacing*{\chapter}{0pt}{0pt}{1.5em}

% Secciones y subsecciones: sobrias y claramente jerarquizadas.
\titleformat{\section}
  {\normalfont\bfseries\fontsize{12}{14.5}\selectfont}
  {\thesection.}
  {0.6em}
  {}
\titlespacing*{\section}{0pt}{1.5em}{0.6em}

\titleformat{\subsection}
  {\normalfont\bfseries\itshape\fontsize{12}{14.5}\selectfont}
  {\thesubsection.}
  {0.6em}
  {}
\titlespacing*{\subsection}{0pt}{1.2em}{0.4em}

\setcounter{secnumdepth}{2}
\setcounter{tocdepth}{2}

% ----------------------------------------------------------------------------
% LISTAS
% ----------------------------------------------------------------------------
\usepackage{enumitem}
\setlist{
  topsep=0.45em,
  itemsep=0.20em,
  parsep=0pt,
  partopsep=0pt
}

% ----------------------------------------------------------------------------
% FIGURAS, TABLAS Y UNIDADES
% ----------------------------------------------------------------------------
\usepackage{graphicx}
\graphicspath{{../figures/}{figures/}}
\usepackage{float}
\usepackage{booktabs}
\usepackage{tabularx}
\usepackage{array}
\usepackage{longtable}
\usepackage{threeparttable}
\usepackage{caption}
\usepackage{subcaption}
\usepackage{siunitx}

\sisetup{
  output-decimal-marker = {,},
  group-separator = {\,},
  group-minimum-digits = 4
}

\captionsetup{
  font=small,
  labelfont=bf,
  labelsep=period,
  justification=justified,
  singlelinecheck=false
}
\captionsetup[table]{position=top}
\captionsetup[figure]{position=bottom}

% Espaciado de flotantes contenido para evitar páginas artificialmente largas.
\setlength{\textfloatsep}{12pt plus 2pt minus 4pt}
\setlength{\floatsep}{10pt plus 2pt minus 3pt}
\setlength{\intextsep}{10pt plus 2pt minus 3pt}
\setlength{\abovecaptionskip}{6pt}
\setlength{\belowcaptionskip}{2pt}

% ----------------------------------------------------------------------------
% CITAS Y REFERENCIAS — APA 7
% ----------------------------------------------------------------------------
\usepackage{csquotes}
\usepackage[
  backend=biber,
  style=apa,
  sorting=nyt,
  giveninits=true,
  uniquename=false,
  doi=true,
  url=true,
  isbn=false
]{biblatex}
\DeclareLanguageMapping{spanish}{spanish-apa}

% Bibliografía canónica de entrega.
\addbibresource{bibliography.bib}

% ----------------------------------------------------------------------------
% HIPERVÍNCULOS Y REFERENCIAS CRUZADAS
% ----------------------------------------------------------------------------
\usepackage{xurl}
\urlstyle{same}
\usepackage[hidelinks]{hyperref}
\usepackage[capitalise,noabbrev]{cleveref}

\hypersetup{
  pdftitle={Differential Machine Learning para la valoración y cobertura de una opción asiática sobre capacidad de cómputo},
  pdfauthor={David Martín Gómez},
  pdfsubject={Máster en Finanzas Cuantitativas},
  pdfkeywords={Differential Machine Learning, opciones asiáticas, RQMC, Log-OU, cobertura, capacidad de cómputo}
}

% Nombres en español homogéneos.
\renewcommand{\contentsname}{Índice general}
\renewcommand{\listfigurename}{Índice de figuras}
\renewcommand{\listtablename}{Índice de tablas}
\addto\captionsspanish{\renewcommand{\tablename}{Tabla}\renewcommand{\listtablename}{Índice de tablas}}

% ----------------------------------------------------------------------------
% DATOS DEL TFM — editar aquí si cambia el título final
% ----------------------------------------------------------------------------
\newcommand{\TFMTitle}{Differential Machine Learning para la valoración y cobertura de una opción asiática sobre capacidad de cómputo}
\newcommand{\TFMSubtitle}{Aplicación al precio por GPU-hora bajo dinámica Log-OU}
\newcommand{\TFMAuthor}{David Martín Gómez}
\newcommand{\TFMProgram}{Máster en Finanzas Cuantitativas}
\newcommand{\TFMInstitution}{UNED — Formación Permanente}
\newcommand{\TFMEdition}{Edición 2025/2026}
\newcommand{\TFMDate}{Septiembre de 2026}

% Input opcional: permite compilar main.tex aunque todavía no existan todos
% los ficheros de capítulos/frontmatter.
\newcommand{\OptionalInput}[1]{%
  \IfFileExists{#1}{\input{#1}}{%
    \typeout{[TFM] Archivo todavía no creado: #1}%
  }%
}

% ----------------------------------------------------------------------------
% DOCUMENTO
% ----------------------------------------------------------------------------
\begin{document}
\pagestyle{plain}

% PORTADA ---------------------------------------------------------------------
\hypersetup{pageanchor=false}
\pagenumbering{roman}
\begin{titlepage}
  \centering
  \thispagestyle{plain}

  \vspace*{1.3cm}

  {\bfseries\fontsize{14}{17}\selectfont \TFMProgram\par}
  \vspace{0.3cm}
  {\fontsize{12}{14.5}\selectfont \TFMInstitution\par}

  \vfill

  \rule{0.86\textwidth}{0.6pt}\par
  \vspace{0.8cm}

  {\bfseries\fontsize{18}{22}\selectfont \TFMTitle\par}
  \vspace{0.55cm}
  {\fontsize{14}{17}\selectfont \TFMSubtitle\par}

  \vspace{0.8cm}
  \rule{0.86\textwidth}{0.6pt}\par

  \vfill

  {\fontsize{12}{14.5}\selectfont
    \textbf{Autor:} \TFMAuthor\par
    \vspace{0.45cm}
    \TFMEdition\par
    \vspace{0.25cm}
    \TFMDate\par
  }

  \vspace*{1.2cm}
\end{titlepage}

% PRELIMINARES ---------------------------------------------------------------
\clearpage
\pagenumbering{roman}
\hypersetup{pageanchor=true}
\setcounter{page}{2}

\chapter*{Resumen}
\addcontentsline{toc}{chapter}{Resumen}
\OptionalInput{frontmatter/resumen.tex}

\clearpage
\chapter*{Abstract}
\addcontentsline{toc}{chapter}{Abstract}
\begin{otherlanguage}{english}
\decimalpoint
\OptionalInput{frontmatter/abstract.tex}
\end{otherlanguage}

\clearpage
\tableofcontents

\clearpage
\listoffigures

\clearpage
\listoftables

% CUERPO PRINCIPAL ------------------------------------------------------------
\clearpage
\pagenumbering{arabic}
\setcounter{page}{1}

\chapter{Introducción}
\label{chap:introduccion}
\OptionalInput{chapters/01_introduccion.tex}

\chapter{Revisión de la literatura y marco conceptual}
\label{chap:literatura}
\OptionalInput{chapters/02_literatura.tex}

\chapter{Datos y mercado de capacidad de cómputo}
\label{chap:datos}
\OptionalInput{chapters/03_datos.tex}

\chapter{Metodología}
\label{chap:metodologia}
\OptionalInput{chapters/04_metodologia.tex}

\chapter{Resultados}
\label{chap:resultados}
\OptionalInput{chapters/05_resultados.tex}

\chapter{Conclusiones}
\label{chap:conclusiones}
\OptionalInput{chapters/06_conclusiones.tex}

% DECLARACIÓN DE IA -----------------------------------------------------------
\IfFileExists{frontmatter/declaracion_ia.tex}{%
  \clearpage
  \chapter*{Declaración de uso de herramientas de inteligencia artificial}
  \addcontentsline{toc}{chapter}{Declaración de uso de herramientas de inteligencia artificial}
  En la elaboración de este trabajo se utilizaron ChatGPT y Codex como herramientas
de apoyo para la revisión de redacción, creación y depuración de código,
comprobaciones de coherencia y revisión de fórmulas y resultados. La revisión final
con Codex se realizó con asistencia basada en GPT-6; no se conserva un registro
completo de las versiones utilizadas en las consultas anteriores de ChatGPT.
Todos los resultados, referencias, fórmulas y procedimientos fueron revisados y
verificados por el autor, que asume íntegramente la responsabilidad sobre el
contenido final.

}{}

% REFERENCIAS -----------------------------------------------------------------
\clearpage
\printbibliography[heading=bibintoc,title={Referencias}]

% ANEXOS ----------------------------------------------------------------------
\IfFileExists{appendices/anexos.tex}{%
  \clearpage
  \appendix
  \input{appendices/anexos.tex}
}{}

\end{document}
